\section{Evaluation}
\label{sec:evaluation}

\subsection{Experimental Framework}
\label{sec:eval-method}

The evaluation quantifies the consequences of the three mechanisms in
Figure~\ref{fig:thesis-map}. For the unified language, it measures model
uncertainty and executable task completion over matched extension tasks. For
mods, it measures the files, source lines, ownership zones, and registration
sites changed by a complete feature. For reactive execution, it measures
edit-to-result latency and work reuse as graph size and affected scope vary.
Separate artifact and overhead measurements establish that responsiveness
preserves correctness, generated-code quality, and acceptable cold-path cost.

\paragraph{Subjects.}
The extension and change-footprint experiments compare Joggle with MLIR and
xDSL. MLIR represents a mature multi-level compiler infrastructure, while xDSL
is a Python-native framework designed for rapid construction of SSA compilers.
Each comparison uses a pinned source revision and the project's documented
extension path. The update experiment uses two tracks. A controlled track
compares Joggle's reactive scheduler with complete
rerun, safe suffix rerun, and one-mechanism-at-a-time ablations inside the same
executable. A cross-system track runs the same analysis and rewrite semantics
over equivalent graph inputs in Joggle, MLIR, and xDSL. The controlled track
identifies the cause of reuse; the cross-system track measures the end-to-end
consequence of each infrastructure's public execution model.

\paragraph{Models.}
The full-model corpus is the repository's pinned set of 16 ONNX models. It
spans classification, detection, machine comprehension, quantized networks,
and vision transformers. Every model passes through a fixed compatibility gate
for each compared system. The corpus record
retains all outcomes; aggregate cross-system results use the common accepted
set. A generated graph family supplements these models at $10^3$, $10^4$,
$10^5$, and $10^6$ operations. It controls total graph size, affected-cone
size, fan-out, and stage count independently.

\paragraph{Correctness.}
Every compiler-function task has a build-and-test oracle. Every transformation
measurement verifies the resulting graph and compares a canonical structural
digest with the expected result. Cross-system outputs use the same semantic
oracle rather than textual equality. Artifact measurements compare output
tensors with a reference implementation under an operator- and dtype-specific
tolerance. A sample enters a latency aggregate only after its correctness
checks pass.

\paragraph{Timing and statistics.}
Systems are compiled in release mode with matched instrumentation.
Cold-process, warm-process, and warm-cache measurements are reported separately.
Each latency case uses ten warm-up iterations followed by 100 timed iterations,
with case order randomized within a block. We report the median, 95th
percentile, and a task- or model-level bootstrap 95\% confidence interval.
Ratios are aggregated with the geometric mean. The artifact records the
machine, operating system, compiler revision and flags, CPU affinity, frequency
policy, model digest, random seed, and raw per-iteration rows.

\subsection{Extension Predictability and Completion}
\label{sec:eval-convenient}

This experiment jointly evaluates predictability and executable task
completion; code length is reported as a control. The suite contains 36 paired
tasks in six families: defining a type or operation, writing an analysis,
applying a graph rewrite, converting a
representation, generating an artifact, and completing a vertical feature
that combines these roles. A task has one semantic specification and a
system-specific harness. Reference solutions use the shortest idiomatic path
that passes the same oracle.

The split is by semantic feature. Definitions, analyses, rewrites, and emitters
for one feature remain in the same partition, preventing a model from seeing
one stage of a held-out feature during context construction. Identifiers,
comments, and formatting follow a fixed normalization policy. Generated files
and test-harness boilerplate are outside the scored continuation.

We evaluate two open-weight code models in the 1--3B parameter range with
frozen weights. Each prompt contains the task specification and a compact API
card. We then vary the number of training-partition demonstrations from zero to
four, selecting them with one deterministic retrieval rule under the same token
budget for every system. Absolute uncertainty captures the complete prompting
cost; the within-task change across demonstration budgets captures how quickly
the extension surface becomes predictable from local examples. Sampling
parameters, seeds, stopping rules, and maximum continuation length are fixed
per model.

For a reference continuation $x_{1:N}$ and context $c$, token perplexity is
\[
  \operatorname{PPL}(x\mid c) =
  \exp\left(-\frac{1}{N}\sum_{t=1}^{N}
  \log p_\theta(x_t\mid x_{<t},c)\right).
\]
We report paired log-perplexity per task and keep values from different
tokenizers separate. Perplexity measures uncertainty over a valid
implementation; pass@$k$ supplies executable validation. A completion must
parse, type-check, compile where required, and pass the task oracle. We report
pass@1, pass@5, and pass@10 from a fixed pool of $n$ samples. If $c$ samples
pass, the estimator is~\cite{chen2021codex}
\[
  \widehat{\operatorname{pass@}k}
    = 1 - \frac{\binom{n-c}{k}}{\binom{n}{k}}.
\]
We also report syntax success, type-check success, context tokens, completion
tokens, and repair-free success. The paired task design separates the extension
surface from differences in task semantics.

\subsection{Change Footprint and Ownership}
\label{sec:eval-controllable}

The change-footprint experiment uses the same semantic feature families but
evaluates repository changes.
Twelve tasks cover new capabilities and revisions to existing capabilities;
each task includes its required semantics, analysis, transformation,
conversion, artifact behavior, and tests. Implementations begin from clean,
pinned snapshots and end only when the common oracle passes. A deterministic
minimization pass then removes each changed hunk in turn and retains the removal
whenever the oracle still passes. This produces one locally minimized,
auditable patch per system and task.

For a patch $p$, we define the primary footprint
\[
  \operatorname{Footprint}(p)=(F_p,L_p,Z_p,R_p),
\]
where $F_p$ is the number of touched implementation files, $L_p$ is added plus
modified implementation lines, $Z_p$ is the number of ownership zones crossed,
and $R_p$ is the number of build, registry, or pipeline declarations changed.
An ownership zone is a source package or build target with a distinct public
responsibility. Test files and test lines are reported in parallel rather than
discarded. Generated files, vendored dependencies, formatter-only changes, and
lock-file churn are excluded by a rule fixed before implementation. Task
specifications, ownership-zone definitions, and counting rules are frozen
before aggregate statistics are computed.

Controllability is evaluated from all four footprint coordinates because line
count and test volume carry different meanings. We show every task and pair the
coordinates with build success and the common semantic oracle. We also report
dependency fan-out: the number of public components whose rebuild or
registration state changes because of the patch. For Joggle, the analysis
records whether the change remains within one mod or crosses declared
\texttt{use} edges.

The principal comparison is paired by task: for each coordinate, we compute
the within-task ratio between Joggle and each baseline, then summarize the
ratios with a geometric mean and bootstrap interval. A second view groups tasks
by role and contrasts small rewrites with vertical features. Complete patches,
counting scripts, and inclusion decisions ship with the artifact.

\subsection{Reactive Update Cost}
\label{sec:eval-efficient}

The update experiment measures an edit-to-result operation: begin with a
verified optimized graph, apply one controlled edit, restore the required
compiler result, and verify its semantic digest. The Joggle-native pipeline
contains analysis,
canonicalization, target selection or legalization, memory planning, and
artifact preparation. The cross-system pipeline uses a common analysis and
local rewrite over graph inputs derived from the same models. Stage order and
rewrite semantics remain fixed across rerun strategies.

The edit matrix covers six invalidation scopes: (1) no graph change, which
isolates validation and cache-hit cost; (2) metadata on one value; (3) the type
of one value; (4) the callee or operands of one operation; (5) one
function-local structural insertion or removal; and (6) a mod dependency or
environment change. Each local edit is instantiated at deterministic positions
selected before timing. Position is blocked by early, middle, and late pipeline
relevance so that one favorable node cannot determine a model's result.
Structural edits preserve verification invariants, and the final oracle checks
that every rerun strategy reaches the same result.

We compare five execution policies. \textbf{Full} runs every stage.
\textbf{Suffix} runs the conservative suffix beginning at the first stage whose
declared input class may have changed. \textbf{Reactive} validates recorded
inputs, propagates overlapping effects, and runs the selected stages.
\textbf{Whole-mod} replaces fine-grained observations with one mod revision.
\textbf{No-plan-cache} retains reactive selection but decodes compiler
functions again. Full and Suffix bound the reusable work. Whole-mod isolates
dependency precision, while No-plan-cache isolates plan reuse.

The primary latency is wall-clock time from the completed edit to a verified
result. Supporting measurements are selection time, evaluation time,
verification time, executed and reused stages, observed functions,
collections, operations, and values, changed functions, compiler-function
operations executed, plan hits, and peak resident memory. We report both
absolute latency and speedup $T_{\mathrm{Full}}/T_p$ for policy $p$. Executed
compiler-function operations $E_p$ define work reuse as
$1-E_p/E_{\mathrm{Full}}$. Work counts explain latency and remain comparable
when host-language runtimes differ.

The 16-model corpus covers realistic graph shapes. The generated family then
separates total graph size $|G|$ from the affected region $|\Delta G|$. One
sweep fixes a one-operation edit and scales $|G|$, exposing the relation
between validation cost and recorded dependency size. A second fixes $|G|$
and increases $|\Delta G|$, locating the crossover at which complete execution
becomes preferable. Every raw sample includes the miss reason and output
digest, which detects unintended reuse.

The cross-system result is reported in two parts. Cold end-to-end execution
starts a fresh process and includes parsing and setup. Warm edit execution keeps
the graph and pass infrastructure resident and applies the same semantic edit
through each system's public API. Cache state is named for every series. This
separation prevents process startup or persistence policy from being
misattributed to dependency-directed execution.

\subsection{Artifact Quality and System Costs}
\label{sec:eval-artifacts}

Artifact measurements separate compiler responsiveness from the code it
produces. The operator
suite covers elementwise chains, reductions, matrix multiplication,
convolution, quantize/dequantize paths, and fusion opportunities over a
predeclared matrix of shapes and dtypes. The model suite uses the compatible
subset of the pinned corpus. Joggle's existing artifact mods instantiate the
same typed generation interface used by the design; generated native sources
are compiled with one fixed toolchain and flag set.

For operators, we report steady-state kernel latency, end-to-end call latency,
generated code size, and compile time. For models, we report per-inference
latency, peak memory, and binary or artifact size. Each result includes the
unoptimized Joggle pipeline, the optimized Joggle pipeline, and a pinned
reference runtime on the same CPU. Speedup is computed per operator or model
before geometric aggregation. Fixed random inputs and framework reference
outputs establish numerical equivalence.

The two comparisons serve distinct purposes. The unoptimized comparison
measures the effect of Joggle extensions at operator and model scale. The
reference-runtime comparison locates the generated artifact relative to an
established implementation. Compilation time remains separate from execution
time, and operator microbenchmarks remain separate from complete-model
measurements.

The final analysis decomposes dependency-directed update cost. A cold run
measures loading, parsing, initial verification, plan construction, and the
first pipeline execution. A warm run measures dependency capture,
dependency validation, transaction setup, graph verification, and commit. Peak
memory is measured after loading, after the first schedule, and after repeated
edits; bytes per recorded function, operation, and value are derived from the
same runs.

Ablations disable dependency capture, fine-grained observations, persistent
plans, dispatch caching, reusable register windows, and lazy structural
snapshots one at a time. Each ablation runs the same inputs and oracles as the
complete system. A final precision study rewrites the same analysis with a
whole-graph traversal, a function-scoped traversal, and direct entity lookup.
It connects API choice to observed dependency size, validation cost, and later
reuse without changing the analysis result.

Together, these experiments form one evidence chain: generation characterizes
the regularity of the extension surface; patch footprint characterizes the
boundary of a complete feature; reactive execution quantifies reuse after an
edit; and artifact and overhead measurements establish the resulting costs.
